\begin{claim}{D.2.41}{Proved}{The actual sphere constraint on fourth-stage relative choices}
Let $M=M(\mathbb Z/2,3)$, let $D=D_2^{(4)}M$ be the
finite quadratic source of D.2.38, and let
$y_0:S^{10}\to C_4M$ be the actual spherical representative
constructed there. Its relative stable class has only a
weight-two component $y_D\in\pi_{10}D$. The fourfold adjunct
of $y_0$ is
\[
 \widetilde y_0=\pm l_4\in\pi_{14}P^8(2),
 \qquad |y_0|=|y_D|=4,
 \tag{D.2.41a}
\]
where $P^8(2)=M(\mathbb Z/2,7)$ and $l_4$ is the
attaching map of the top cell of the Stiefel manifold $V_{9,2}$.

Put $Y=\Sigma^2K(\mathbb F_2,2)$,
$Z'=Y\langle7\rangle$, $U=\Omega^4Z'$, and
$V=U\langle6\rangle$. Thus $V$ is five-connected.
Let $m:M\to U$ be the canonical dangerous Moore restriction
of D.2.29, and let $g_4:C_4M\to U$ be its free $E_4$
extension. If an ordinary map $f:QM\to U$ restricts to $m$,
its difference from $g_4$ determines a weight-two relative map
$v_2:D\to V$ satisfying
\[
 (v_2)_*y_D=-(g_4)_*y_0\in2\pi_{10}V,
 \qquad (v_2)_*Y_{10}=0\in H_{10}(V;\mathbb F_2).
 \tag{D.2.41b}
\]
The identification $\pi_{10}V=\pi_{10}U$ is used in this formula.
In particular, for the fixed paired functional
\[
 \mu_{11}=\sigma^{*2}
 \bigl((\xi_1k)^*+(\xi_1f)^*\bigr)\in H^{11}U
 \tag{D.2.41c}
\]
of D.2.40, the target relation
\[
 \left.\mu_{11}\right|_V\in
 \operatorname{Sq}^1H^{10}V+
 \operatorname{Sq}^3H^8V+
 \operatorname{Sq}^5H^6V
 \tag{D.2.41d}
\]
would suffice to exclude such a full extension. All cohomology
in \textup{(D.2.41c)--(D.2.41d)} has coefficients $\mathbb F_2$.
Relation \textup{(D.2.41d)} is a sufficient test, not a proved
target identity. The fourth-suspension extension problem and
Part D.2.4 remain open.
\end{claim}
\begin{proof}
\textbf{The arity-three component.}
Write $A_0,B_0$ for the bottom and top homology classes of $M$.
The free $E_4$ homology basis gives the following complete list
through degree eleven in $E_3=D_3^{(4)}M$:
\[
\begin{array}{c|l}
9&A_0^3\\
10&A_0^2B_0,\ A_0Q^4A_0\\
11&A_0B_0^2,\ A_0Q^5A_0,\ B_0Q^4A_0.
\end{array}
 \tag{D.2.41e}
\]
An indecomposable Dyer--Lashof expression has weight a power
of two. In this range the weight-three classes are therefore
products of weights one and two: the first weight-three
Browder expression has degree at least twelve. This proves
the completeness of the list, using
\cite[Chapter III, Theorems~3.1--3.2]{I_CohenLadaMay1976}.
Cartan and the operation table in D.2.38 give
\[
\begin{split}
 \beta(A_0^2B_0)&=\beta(A_0Q^4A_0)=A_0^3,\\
 \beta(B_0Q^4A_0)&=A_0^2B_0+A_0Q^4A_0,
\end{split}
 \tag{D.2.41f}
\]
with the other two degree-eleven Bocksteins zero. Each of the
three degree-eleven classes has second square $A_0^3$.

Consequently the stable ten-skeleton is
$M(\mathbb Z/2,9)\vee S^{10}_c$, where $c$ reduces to
$A_0^2B_0+A_0Q^4A_0$. Normalize the nonzero eleven-cell
integral boundary to $2c$ by an odd $2$-local unit. The other
two boundaries are divisible by four and can be made zero by
even cell slides. These slides preserve all the mod-two square
vectors. The cell $A_0B_0^2$ then attaches by the sole nonzero
element $i\eta_9\in\pi_{10}M(\mathbb Z/2,9)$, detected by
its second square. It kills that torsion summand; the remaining
relation is $2c$. Thus
\[
 \pi_{10}E_3=\mathbb Z/2,
 \qquad
 h_2:\pi_{10}E_3\longrightarrow H_{10}(E_3;\mathbb F_2)
 \text{ is injective}.
 \tag{D.2.41g}
\]
These are stable cell models, not an unstable cell cancellation
claim. Since $E_3$ is eight-connected, its ordinary group in
this degree agrees with the stable group. Cells above degree
eleven do not affect it. The degree-two boundary normalization
also follows directly from the rank-one first Bockstein in
\textup{(D.2.41f)}; there is no assumption that a nonzero
boundary has an unspecified higher power of two.

The relative Snaith splitting of $C_4M/M$ is compatible with
weights; see \cite[Theorems~1.1 and~3.2]{Snaith1974}.
The actual sphere $y_0$ has homology solely in the weight-two
class $Y_{10}$. Its weight-three projection is therefore zero
by \textup{(D.2.41g)}. The summands of weight at least four
are eleven-connected and contribute nothing to $\pi_{10}$.
Its entire relative stable class is consequently its
weight-two projection $y_D$. D.2.38 gives $|y_D|=4$.

\medskip\noindent
\textbf{The actual sphere before stable projection.}
D.2.38 proves that $y_0$ has zero image in $QM$. Thus its
fourfold adjunct $\widetilde y_0:S^{14}\to P^8(2)$ has zero
stable image. The primary calculation
\cite[Lemmas~2.7 and~2.8(ii), p.~523]{EH5MorisugiMukai2000}
is
\[
\begin{split}
 \pi_{14}P^8(2)&=\mathbb Z/4\{l_4\}
       \oplus\mathbb Z/2\{\widetilde{\nu_7^2}\},\\
 2l_4&=i_8\sigma',\qquad
 \Sigma l_4=i_9[\iota_8,\iota_8].
\end{split}
 \tag{D.2.41h}
\]
The second summand pinches to $\nu_8^2$, whose stable image is
nonzero. The $l_4$ summand stabilizes to zero by the displayed
Whitehead-square formula. Hence the exact stable kernel is
$\mathbb Z/4\{l_4\}$. Since the relative projection of
$\widetilde y_0$ has order four, its coefficient on $l_4$ is
odd. This proves \textup{(D.2.41a)}, including the order of
the actual sphere rather than merely its stable projection.

\medskip\noindent
\textbf{The relative map and its precise sphere value.}
The space $C_4M/M$ is five-connected. Form the pointwise
difference of $f|_{C_4M}$ and $g_4$ in the H-group $U$,
using the chosen homotopy between their restrictions to $M$.
The cofiber universal property produces a map
$d:C_4M/M\to U$, and connectivity gives its unique lift
to $V$. No multiplicativity of $f$ is used.

For a finite source $K$ of dimension at most eleven, stabilization
gives a bijection
\[
 [K,V]\ \xrightarrow{\ \cong\ }\
 [\Sigma^\infty K,\Sigma^\infty V].
 \tag{D.2.41i}
\]
Indeed the five-connectivity of $V$ makes the first suspension
surjective in this range. The unit $V\to\Omega\Sigma V$
has an actual left inverse because $V$ is a loop space, which
gives injectivity. The space $\Sigma V$ is six-connected and
$\dim\Sigma K\leq12$, so every subsequent suspension is
bijective by Freudenthal. Apply \textup{(D.2.41i)} to the
weight-two stable restriction of $d$ to obtain $v_2:D\to V$.
Apply it also to $S^{10}$: the weight calculation above and
the nullity of $y_0$ in $QM$ give the actual equality
\[
 (v_2)_*y_D=d_*y_0=-(g_4)_*y_0.
 \tag{D.2.41j}
\]
The middle notation means composition with the quotient
$C_4M\to C_4M/M$. This equality is proved in $\pi_{10}V$
and then identified with $\pi_{10}U$; it does not require
injectivity of the map of suspension spectra induced by
$V\to U$.

Write $\alpha\in\pi_3U=\mathbb Z/4$ and
$\rho\in\pi_4U=\mathbb Z/4$ for the generators already
computed, so $\alpha\eta_3=2\rho$. The bottom restriction
of $m$ is $-2\alpha$. In the stable range of maps $M\to M$,
the Moore Puppe sequence factors $2\id_M$ through the top
pinch followed by either zero or $i\eta_3$. Therefore
\[
 2m=0,
 \qquad
 \lambda_3(-2\alpha,m)
   =-2\lambda_3(\alpha,m)
   =-\lambda_3(\alpha,2m)=0.
 \tag{D.2.41k}
\]
For the first equality, the possibly nonzero factor gives
$(-2\alpha)\eta_3q=-4\rho q=0$. The second uses genuine
bi-additivity: the desuspensions of the two inputs of the
fourfold-adjoint Whitehead product are $S^6$ and $\Sigma^3M$,
both suspensions. Thus the hypotheses of
\cite[Proposition~3.4, p.~14]{EH5Arkowitz1962} apply.

The Moore map whose top-sphere factor is $(g_4)_*y_0$ is
exactly the generalized Browder product in
\textup{(D.2.41k)}. The actual Puppe sequence for
$M(\mathbb Z/2,9)$ now implies
$(g_4)_*y_0\in2\pi_{10}U$. Since $\pi_{10}V\to\pi_{10}U$
is an isomorphism, its lift is also a double. Equations
\textup{(D.2.41j)--(D.2.41k)} prove \textup{(D.2.41b)}.
They do not assert this doubled element is zero. More precisely,
if $h_M:P^8(2)\to Y$ is the adjoint of $m$, its value is
\[
 (g_4)_*y_0=\pm h_Ml_4\in\pi_{14}Y,
 \qquad 2(h_Ml_4)=(-2\alpha)\sigma'.
 \tag{D.2.41l}
\]
In the final formula $\alpha$ denotes its corresponding
generator of $\pi_7Y$. The composite $h_Ml_4$ is not evaluated.

\medskip\noindent
\textbf{The weaker sufficient target identity.}
The source table in D.2.38 gives
\[
 \beta Z_{11}=Y_{10},\qquad
 \operatorname{Sq}_*^3Z_{11}
   =\operatorname{Sq}_*^5Z_{11}=0.
 \tag{D.2.41m}
\]
Consequently \textup{(D.2.41b)} makes every term on the
right of \textup{(D.2.41d)} evaluate to zero on
$(v_2)_*Z_{11}$. The first-square term uses the specific
sphere constraint; it need not vanish for an arbitrary map
$D\to V$.

The top class $Q^7B_0$ is primitive, so the pointwise difference
adds its two homology images without a convolution term.
For the canonical $g_4$ its $\mu_{11}$ value is one, by
D.2.29 and D.2.40. For a map from $QM$ the value is zero:
$Q^7B_0=\beta Q^8B_0$ in the infinite-loop source, whereas
neither outer restricted-square coefficient defining
\textup{(D.2.41c)} occurs in $\beta H_{14}T$. The latter is
the free restricted-Lie basis argument of D.2.29 and applies
to both $\xi_1k$ and $\xi_1f$. Thus a full extension would
require the relative value one, contradicting
\textup{(D.2.41d)} if that relation were established.

Finally, D.2.36 and the Moore Puppe sequence show that every
dangerous fourth-suspension map has this same Moore restriction:
for fixed bottom map, the only nontrivial top change is
$\rho q$, whose double homology suspension adds
$r=b^2+k+p_3$. The entire-unit linear exclusion of D.2.36
rules out that change. Hence the sufficient test has the
stated relevance to all dangerous fourth-suspension maps,
without fixing their entire infinite unit maps. D.2.42 below
disproves the target membership in this conditional test.
\end{proof}
